\section{Introduction}
\label{sec:introduction}

\subsection{Research context}

This project investigates a computational workflow for urban flood forecasting
that combines a numerical shallow-water model, a GPU-oriented implementation
of the same governing physics, and a learned residual correction.  The
research objective is not to replace the governing solver with a purely
data-driven surrogate.  Instead, the architecture assigns distinct roles to
three components: ANUGA provides a numerical reference solution, a custom JAX
finite-volume solver provides the accelerated forecast trajectory, and a
residual network learns systematic differences between their outputs.

The emphasis is therefore on the interaction between physical modelling,
numerical verification, and computational performance.  The repository is
organised so that raw GIS and rainfall inputs remain read-only, derived
datasets are reproducible, and report figures are stored with the generated
model outputs.  Configuration is composed through Hydra-style component files,
with a project-wide random seed of 42 and a common projected coordinate
reference system of EPSG:32643.

\subsection{Aim and objectives}

The aim is to assess whether a JAX-based physics--AI workflow can support
faster forward simulation while retaining agreement with the ANUGA numerical
reference.  Within the implemented research design, this aim is
decomposed into the following objectives:

\begin{enumerate}
    \item validate and reproject the available boundary, mesh, roughness, and
    rainfall inputs without modifying the raw source files;
    \item construct a deterministic synthetic terrain scenario because no
    surveyed DEM is currently available;
    \item establish an ANUGA numerical reference and an independently
    implemented structured-grid JAX shallow-water solver;
    \item test water conservation and preservation of a lake-at-rest state;
    \item train a residual model on chronologically separated solver-pair data
    and couple its corrections back into the JAX time integration; and
    \item compare matched JAX and PyTorch residual networks under a controlled
    GPU benchmarking protocol.
\end{enumerate}

\subsection{Scope and interpretation}

The present repository supports a controlled methodological study rather than
an operational forecast for Gurugram.  ANUGA is treated as a numerical
reference, not observational ground truth.  No observed flood depths,
discharges, surveyed elevations, or calibration measurements are currently
included.  The available machine-learning dataset comprises 60 chronological
transitions from one simulated event, so its held-out time block does not
demonstrate generalisation to independent storms.

Two further restrictions define the current evidential boundary.  First, the
meaning of the source boundary attributes \texttt{BndTypeNo} and
\texttt{ConstValue} could not be recovered from the available files.  A final
methodological decision therefore treats every boundary as reflective.  This
is a deliberate closed-domain simplification and is reported as a limitation,
not as an interpretation of the source schema.  Second, the configured
development simulation covers one hour of a 35-hour rainfall record and
excludes the event peak.  Results from that window are suitable for testing
the end-to-end pipeline, but not for characterising a complete return-period
event.

\subsection{Repository-based reproducibility}

The reusable Python package and executable entry points are maintained
separately under the source and scripts directories.  Raw inputs are separated
from interim, synthetic, processed, and output products.  Numerical tests and
the LaTeX report are also maintained in dedicated directories.  Python 3.11
dependencies and optional benchmark dependencies are managed through
\texttt{uv}.

CUDA-enabled JAX packages are part of the environment, but actual GPU
execution requires an exposed NVIDIA device and a working host driver.  The
runtime can fall back explicitly to CPU for development, whereas report-quality
performance measurements are configured to reject such fallback.  This
separation prevents CPU timings from being presented as GPU evidence.
